% Appendix: beating FCFS without observing the FCFS queue.
% Compiles standalone. To include in the paper, copy everything between
% "BEGIN APPENDIX BODY" and "END APPENDIX BODY" (the body uses only amsmath,
% amssymb, amsthm and the macros and theorem environments declared below).
\documentclass[11pt]{article}
\usepackage[margin=1.2in]{geometry}
\usepackage{amsmath,amssymb,amsthm}

\newtheorem{bthm}{Theorem}[section]
\newtheorem{blem}[bthm]{Lemma}
\newtheorem{bprop}[bthm]{Proposition}
\theoremstyle{definition}
\newtheorem{bdef}[bthm]{Definition}
\theoremstyle{remark}
\newtheorem{brem}[bthm]{Remark}

\newcommand{\E}{\mathbb{E}}
\renewcommand{\Pr}{\mathbb{P}}
\newcommand{\N}{\mathbb{N}}
\newcommand{\ind}{\mathbf{1}}
\newcommand{\cX}{\mathcal{X}}
\newcommand{\cF}{\mathcal{F}}
\newcommand{\cL}{\mathcal{L}}
\newcommand{\Reps}{R^{\varepsilon}}
\newcommand{\gp}{\mathsf{g}}
\newcommand{\SvA}{\Sigma_A}
\newcommand{\SvB}{\Sigma_B}
\newcommand{\Poi}{\mathrm{Poi}}

\begin{document}
\appendix
\renewcommand{\thesection}{B}

% ============================ BEGIN APPENDIX BODY ============================
\section{Beating FCFS without observing its state}\label{app:blind}

The Ledger rule conditions on the state of the FCFS queue. This appendix shows
that this is not necessary: queue $A$ can attract a strict majority of arrivals
with a service rule that depends only on $A$'s own history. The rule uses two
inferences. First, if an agent joins $A$ in slot $2$, then queue $B$ held at
least two agents at that moment. Second, if $A$ then completes three services
in a very short time, queue $B$ has almost surely not completed any, so it
still holds at least two agents. Queue $A$ is then empty and can offer slot $1$
rather than slot $0$: the offer still attracts the next arrival, and slot $0$
is held in reserve for the arrival after that.

Throughout, the model and notation are those of Sections~3 and~5.1: slots are
indexed by $\N_0$, the service rate is $1$ at each queue, arrivals have rate
$\lambda\in(0,2)$, queue $B$ uses FCFS, and agents minimize their expected
waiting time, breaking ties with the draws $U_n$. We write $a(t)=|S_A(t)|$ and
$b(t)=|S_B(t)|$, and all processes are taken right-continuous with left limits.
Both queues start empty.

\subsection{The rule and the result}

\begin{bdef}\label{def:blind}
A service rule for queue $A$ is \emph{blind} if its three mappings depend on
the history $N(t)$ only through the current time, the service times of queue
$A$, and the times at which agents joined queue $A$ together with the slots
they were assigned.
\end{bdef}

A blind rule cannot condition on $S_B$, on arrivals that route to $B$, or on
the service times of $B$.

\begin{bdef}[The rule $\Reps$]\label{def:Reps}
Fix $\varepsilon\in(0,1]$. At time $t$, let $\ell(t)$ be the time of the most
recent join to queue $A$ strictly before $t$ and let $j(t)$ be the slot
assigned at that join (if no agent has joined $A$, set $j(t)=-1$). Say that
queue $A$ is \emph{armed} at $t$, and write $D(t)=1$, if
\[
  S_A(t)=\emptyset,\qquad j(t)=2,\qquad t-\ell(t)<\varepsilon ;
\]
otherwise $D(t)=0$. The rule $\Reps$ is:
\begin{enumerate}
\item \emph{Placement.} An agent arriving at time $t$ is offered slot $1$ if
  $D(t^-)=1$, and slot $\min(\N_0\setminus S_A(t^-))$ otherwise.
\item \emph{Service selection.} $\sigma_A=\min S_A$.
\item \emph{Reassignment.} $\rho_A(i)=i-1$ for all $i\in S_A\setminus\{\sigma_A\}$.
\end{enumerate}
\end{bdef}

The rule $\Reps$ is blind and work-conserving. It differs from FCFS only while
armed, that is, only when $A$ is empty, its last arrival was placed in slot $2$,
and this happened less than $\varepsilon$ ago (so that $A$ has completed three
services within $\varepsilon$).

\begin{bthm}\label{thm:blind}
Fix $\lambda\in(0,2)$. There exist constants $\bar\varepsilon\in(0,1]$ and
$\kappa>0$, depending only on $\lambda$, such that for every
$\varepsilon\in(0,\bar\varepsilon)$ the following holds. If queue $A$ uses the
blind rule $\Reps$, queue $B$ uses FCFS, and agents route optimally, then the
asymptotic arrival rates $\alpha_A=\lim_{t\to\infty}N_A(t)/t$ and
$\alpha_B=\lim_{t\to\infty}N_B(t)/t$ exist almost surely and
\[
  \alpha_A\;\ge\;\frac{\lambda}{2}
  +\frac{\lambda^{2}\varepsilon^{4}}{2}
  \Bigl(\kappa-\frac{\varepsilon^{2}}{12(1+\lambda)}\Bigr)\;>\;\frac{\lambda}{2}.
\]
\end{bthm}

The remainder of this appendix proves Theorem~\ref{thm:blind}. The constants
are $\kappa=\delta_2\,q/(384\,e^{4})$ and
$\bar\varepsilon=\min\{1,\sqrt{12(1+\lambda)\kappa}\}$, where $\delta_2$ is
defined in Lemma~\ref{lem:gap} and $q$ in Lemma~\ref{lem:occupation}.

\subsection{Reachable states and optimal routing}

Write $[k]=\{0,1,\dots,k-1\}$ for $k\ge 0$, so that $[0]=\emptyset$.

\begin{blem}\label{lem:states}
Under $\Reps$, at every time $S_A(t)\in\{[k]:k\ge0\}\cup\{\{1\}\}$, whatever
the routing decisions of the agents.
\end{blem}

\begin{proof}
The claim holds at $t=0$. Suppose $S_A=[k]$. A service completion gives
$[k-1]$ if $k\ge1$ and leaves $S_A=\emptyset$ if $k=0$. An agent who joins is
placed in slot $k$, giving $[k+1]$, unless $A$ is armed, in which case $k=0$
and she is placed in slot $1$, giving $\{1\}$. Suppose $S_A=\{1\}$. Then $A$ is
not armed, so an agent who joins is placed in slot $0$, giving $[2]$, and a
service completion serves the agent in slot $1$, giving $\emptyset$.
\end{proof}

We therefore describe the system by the pair $X(t)=(x_A(t),b(t))$ together
with the indicator $D(t)$, where $x_A=k$ if $S_A=[k]$ and $x_A=\gp$ if
$S_A=\{1\}$. The state space of $X$ is
\[
  \cX=\{(k,b):k,b\in\N_0\}\cup\{(\gp,b):b\in\N_0\}.
\]
We call states $(k,b)$ \emph{gap-free} and states $(\gp,b)$ \emph{gap states}.
For $x\in\cX$ let $a(x)$ and $b(x)$ denote the number of agents in $A$ and in
$B$ (so $a(\gp,b)=1$), and let $n(x)=a(x)+b(x)$.

The waiting time of an agent is the time from her arrival to her own service
completion, so that an agent who joins the FCFS queue $B$ when it holds $b$
agents has expected waiting time $b+1$.

\begin{blem}[Optimal routing]\label{lem:routing}
Suppose $A$ uses $\Reps$ and $B$ uses FCFS. Consider an agent who arrives when
the state is $x$ and the armed indicator is $d$. Her optimal choice is the
following, and does not depend on the strategies of other agents except in
case~\emph{(iii)}, where it is optimal given that later agents route optimally.
\begin{enumerate}
\item[(i)] If $x=(k,b)$ and $d=0$: she joins $A$ if $k<b$, joins $B$ if $k>b$,
  and is indifferent if $k=b$.
\item[(ii)] If $x=(\gp,b)$: she joins $A$ if $b\ge1$ and is indifferent if
  $b=0$.
\item[(iii)] If $x=(0,b)$ and $d=1$: she joins $A$ if $b\ge1$ and joins $B$ if
  $b=0$.
\end{enumerate}
\end{blem}

\begin{proof}
Joining $B$ gives expected waiting time $b+1$ in every case.

(i) At $A$ she is placed in slot $k$ and slots $0,\dots,k-1$ are occupied.
Since agents advance one slot at each service completion and no slot below her
is ever empty, no later arrival can be placed below her, and she is served at
exactly the $(k+1)$st service time of $A$ after her arrival. Her expected
waiting time at $A$ is $k+1$.

(ii) At $A$ she is placed in slot $0$ and is served at the next service time of
$A$. Her expected waiting time at $A$ is $1$, against $b+1$ at $B$.

(iii) At $A$ she is placed in slot $1$, and $S_A$ becomes $\{1\}$. Let $s$ be
the first service time of $A$ after her arrival, and let $\pi_b$ be the
probability that some later agent joins $A$ before $s$. Until $s$ or such a
join, $S_A=\{1\}$, and by (ii) a later agent's decision in this situation does
not depend on anyone else's strategy. If no agent joins before $s$, she is
served at $s$. If an agent joins before $s$, that agent is placed in slot $0$
and is served at $s$, after which our agent occupies slot $0$ and is served at
the following service time of $A$. Since the event that an agent joins before
$s$ is determined by time $s$ and the next service of $A$ takes a further
expected time $1$, her expected waiting time at $A$ is $1+\pi_b$.

Any agent who joins before $s$ arrives before $s$, so
$\pi_b\le\lambda/(1+\lambda)<1$. Hence if $b\ge1$ her expected waiting time at
$A$ is strictly below $2\le b+1$, and she joins $A$. If $b=0$, then with
probability $\lambda/(1+\lambda)$ the next arrival occurs before $s$; at that
moment $B$ is still empty, the arriving agent is indifferent by (ii), and she
joins $A$ with probability $1/2$. Thus $\pi_0\ge\lambda/(2(1+\lambda))>0$, the
expected waiting time at $A$ is strictly above $1$, and she joins $B$.
\end{proof}

\begin{brem}
Case~(i) with $k=2$ gives the first inference: an agent who joins $A$ in slot
$2$ reveals that $b\ge 2$ at that moment.
\end{brem}

\paragraph{Transitions.}
By Lemmas~\ref{lem:states} and~\ref{lem:routing}, the process $X$ moves as
follows. At a service time of $A$: $(k,b)\to((k-1)^+,b)$ and $(\gp,b)\to(0,b)$.
At a service time of $B$: $(k,b)\to(k,(b-1)^+)$ and $(\gp,b)\to(\gp,(b-1)^+)$.
At an arrival with tie-breaking draw $u\in\{A,B\}$ and armed indicator
$d\in\{0,1\}$ immediately before the arrival:
\begin{align*}
d=0:\quad & (k,b)\to
  \begin{cases}(k+1,b)&\text{if }k<b,\text{ or }k=b\text{ and }u=A,\\
               (k,b+1)&\text{if }k>b,\text{ or }k=b\text{ and }u=B;\end{cases}\\
& (\gp,b)\to
  \begin{cases}(2,b)&\text{if }b\ge1,\text{ or }b=0\text{ and }u=A,\\
               (\gp,1)&\text{if }b=0\text{ and }u=B;\end{cases}\\
d=1:\quad & (0,b)\to
  \begin{cases}(\gp,b)&\text{if }b\ge1,\\ (0,1)&\text{if }b=0.\end{cases}
\end{align*}
For $f:\cX\to\mathbb{R}$ and $d\in\{0,1\}$ define the corresponding generator
\[
  \cL_d f(x)=\lambda\bigl(\E_u f(\mathrm{arr}_d(x,u))-f(x)\bigr)
  +\bigl(f(\mathrm{svc}_A(x))-f(x)\bigr)+\bigl(f(\mathrm{svc}_B(x))-f(x)\bigr),
\]
where $\mathrm{arr}_d$, $\mathrm{svc}_A$, $\mathrm{svc}_B$ are the three maps
just described, $u$ is uniform on $\{A,B\}$, and $\cL_1$ is only applied at
states $(0,b)$. The generator $\cL_0$ is defined on all of $\cX$; it is the
generator of the Markov chain $X^0$ on $\cX$ in which $A$ always offers its
lowest empty slot. Restricted to gap-free states, $X^0$ is the chain induced by
two FCFS queues with shortest-queue routing.

Let $\SvA$ and $\SvB$ denote the Poisson processes of service times, $N$ the
Poisson process of arrivals, and $(\cF_t)$ the filtration generated by these
processes and the tie-breaking draws of agents who have arrived. For $s<t$ we
write $\SvA(s,t]$ for the number of service times of $A$ in $(s,t]$, and
similarly for $\SvB$ and $N$.

\begin{blem}[Compensation]\label{lem:dynkin}
Let $f:\cX\to\mathbb{R}$ satisfy $|f(x)|\le c\,(1+n(x))^{2}$ for some constant
$c$. Under $\Reps$ with optimal routing,
\[
  M^f(t)=f(X(t))-f(X(0))-\int_0^t \cL_{D(s)}f(X(s))\,ds
\]
is an $(\cF_t)$-martingale. The same holds for the chain $X^0$ with $\cL_0$ in
place of $\cL_{D(s)}$, from any initial state.
\end{blem}

\begin{proof}
The process $X$ is piecewise constant and jumps only at points of $\SvA$,
$\SvB$ and $N$, and the jump at such a point is the image of $X(s^-)$ under
$\mathrm{svc}_A$, $\mathrm{svc}_B$ or $\mathrm{arr}_{D(s^-)}(\cdot,u)$
respectively, where $u$ is the draw of the arriving agent. Hence
$f(X(t))-f(X(0))$ is a sum of integrals of the predictable integrands
$f(\mathrm{svc}_A(X(s^-)))-f(X(s^-))$, $f(\mathrm{svc}_B(X(s^-)))-f(X(s^-))$ and
$f(\mathrm{arr}_{D(s^-)}(X(s^-),u))-f(X(s^-))$ against $\SvA$, $\SvB$ and the
marked Poisson process of arrivals. Since $n(X(s))\le n(X(0))+N(s)$, each
integrand is bounded in absolute value by $2c\,(2+n(X(0))+N(t))^{2}$ on
$[0,t]$, which has finite mean, and the number of jumps on $[0,t]$ has finite second moment. The
compensation formula for marked Poisson processes (see, e.g.,
Br\'emaud, 1981, Chapter~VIII) therefore shows that subtracting
$\int_0^t\cL_{D(s^-)}f(X(s^-))\,ds$ yields a martingale. Finally $X(s^-)=X(s)$
and $D(s^-)=D(s)$ for all but countably many $s$, so the integral equals
$\int_0^t\cL_{D(s)}f(X(s))\,ds$.
\end{proof}

By the strong Markov property of the driving Poisson processes, the same
conclusion holds for the process restarted at any stopping time, and the
optional stopping theorem applies at bounded stopping times.

\subsection{Reduction to busy times}

Let $r:\cX\to\{-1,0,1\}$ be
\[
  r(x)=\ind\{a(x)>0\}-\ind\{b(x)>0\}.
\]

\begin{blem}\label{lem:busy}
Under $\Reps$ with optimal routing, for every $T\ge0$,
\[
  \E\bigl[N_A(T)-N_B(T)\bigr]
  =\E\int_0^T r(X(s))\,ds+\E\bigl[a(X(T))-b(X(T))\bigr].
\]
\end{blem}

\begin{proof}
Let $C_q(T)$ be the number of service completions at queue $q$ up to time $T$.
Every agent who has joined $q$ has either been served or is still in $q$, so
$N_A(T)=C_A(T)+a(X(T))$ and $N_B(T)=C_B(T)+b(X(T))$. A service time of $A$ is a
service completion exactly when $A$ is non-empty immediately before it, so
$C_A(T)=\int_0^T\ind\{a(X(s^-))>0\}\,d\SvA(s)$ and
$\E\,C_A(T)=\E\int_0^T\ind\{a(X(s))>0\}\,ds$. The same holds for $B$.
\end{proof}

Thus $A$ attracts more arrivals than $B$ in the long run exactly when it is
busy for a larger fraction of time.

\subsection{The value function of FCFS}

We now construct a function $h$ on $\cX$ that solves the Poisson equation
$\cL_0h=-r$. On gap-free states, $h(k,b)$ measures the long-run advantage, in
expected service completions, that $A$ derives from starting in $(k,b)$ when
both queues offer their lowest empty slot.

For the chain $X^0$ started from a gap-free state $(k,b)$, let
$\tau=\inf\{t\ge0:a(X^0(t))=b(X^0(t))\}$. Define
\begin{equation}\label{eq:hdef}
  h(k,b)=
  \begin{cases}
    \E_{(k,b)}\displaystyle\int_0^{\tau}\ind\{b(X^0(s))=0\}\,ds & \text{if }k>b,\\[1ex]
    0 & \text{if }k=b,\\[0.5ex]
    -h(b,k) & \text{if }k<b.
  \end{cases}
\end{equation}

\begin{blem}\label{lem:h}
The function $h$ in \eqref{eq:hdef} is well defined and satisfies:
\begin{enumerate}
\item[(a)] $0<h(k,b)\le (k-b)/\lambda$ for all $k>b$; hence
  $|h(k,b)|\le |k-b|/\lambda$ for all $k,b$.
\item[(b)] $h(1,0)=1/(1+\lambda)$.
\item[(c)] $\cL_0h(x)=-r(x)$ for every gap-free state $x$.
\end{enumerate}
\end{blem}

\begin{proof}
(a) Let $k>b$. On $[0,\tau)$ we have $a>b\ge0$, so every arrival joins $B$ and
$A$ is busy. Applying Lemma~\ref{lem:dynkin} to $f(x)=a(x)-b(x)$, for which
$\cL_0f(x)=-\lambda-1+\ind\{b(x)>0\}\le-\lambda$ whenever $a(x)>b(x)$, and
stopping at $\tau\wedge t$, gives
\[
  0\le\E_{(k,b)}\bigl[a(X^0(\tau\wedge t))-b(X^0(\tau\wedge t))\bigr]
  \le (k-b)-\lambda\,\E_{(k,b)}[\tau\wedge t].
\]
Letting $t\to\infty$ gives $\E_{(k,b)}\tau\le(k-b)/\lambda$, so $h$ is well
defined and $h(k,b)\le\E_{(k,b)}\tau\le(k-b)/\lambda$. For strict positivity,
with probability $(2+\lambda)^{-b}$ the first $b$ events are service times of
$B$, after which the chain is in state $(k,0)$ with $k\ge1$ and remains there
for an expected time of $1/(1+\lambda)$ before $\tau$ can occur. Hence
$h(k,b)\ge(2+\lambda)^{-b}/(1+\lambda)>0$.

(b) From $(1,0)$ the chain stays put until the first arrival or service time
of $A$, at which point it moves to $(1,1)$ or $(0,0)$. So $\tau$ is exponential
with rate $1+\lambda$ and $b=0$ throughout $[0,\tau)$.

(c) Let $x=(k,b)$ with $k>b$. Then $r(x)=\ind\{b=0\}$. Conditioning on the
first jump of $X^0$ and using the strong Markov property, together with
$h=0$ on the diagonal, gives
\[
  h(x)=\frac{1}{\Lambda(x)}\Bigl(\ind\{b=0\}+\sum_{y\ne x}Q_0(x,y)\,h(y)\Bigr),
\]
where $Q_0$ denotes the jump rates of $X^0$ and
$\Lambda(x)=\sum_{y\ne x}Q_0(x,y)$. Rearranging, $\cL_0h(x)=-\ind\{b=0\}=-r(x)$.
Next let $\iota(k,b)=(b,k)$. The dynamics of $X^0$ on gap-free states are
symmetric under exchanging the roles of the two queues, so
$\cL_0(f\circ\iota)=(\cL_0f)\circ\iota$ for every $f$. Since $h\circ\iota=-h$
and $r\circ\iota=-r$, for $k<b$ we obtain
$\cL_0h(k,b)=-(\cL_0h)(b,k)=r(b,k)=-r(k,b)$. Finally, on the diagonal,
\[
  \cL_0h(k,k)=\tfrac{\lambda}{2}\bigl(h(k+1,k)+h(k,k+1)\bigr)
  +\ind\{k>0\}\bigl(h(k-1,k)+h(k,k-1)\bigr)=0=-r(k,k).\qedhere
\]
\end{proof}

We extend $h$ to gap states. For $X^0$ started from $(\gp,b)$, let $\zeta$ be
the first time at which $X^0$ is gap-free, and define
\begin{equation}\label{eq:hgap}
  h(\gp,b)=\E_{(\gp,b)}\Bigl[\int_0^{\zeta}r(X^0(s))\,ds+h(X^0(\zeta))\Bigr].
\end{equation}
The time $\zeta$ is at most the first service time of $A$, so $\E\zeta\le1$,
and $|h(X^0(\zeta))|\le(2+b+N(\zeta))/\lambda$ by Lemma~\ref{lem:h}(a); thus
\eqref{eq:hgap} is well defined. Conditioning on the first jump as in the proof
of Lemma~\ref{lem:h}(c) gives $\cL_0h(\gp,b)=-r(\gp,b)$. Therefore
\begin{equation}\label{eq:poisson}
  \cL_0h(x)=-r(x)\qquad\text{for all }x\in\cX .
\end{equation}

\subsection{The value of a gap}

The next lemma is the heart of the argument. It says that, holding the number
of agents in each queue fixed, it is strictly better for $A$ to hold its single
agent in slot $1$ than in slot $0$.

\begin{blem}\label{lem:gap}
For $b\ge0$ let $\delta_b=h(\gp,b)-h(1,b)$. Then
\[
  0<\delta_b\le\bar\delta:=\max\{2h(2,1),\,h(2,0)\}\qquad\text{for all }b\ge0.
\]
Consequently $|h(x)|\le n(x)/\lambda+\bar\delta$ for all $x\in\cX$.
\end{blem}

\begin{proof}
Fix $b$ and run two copies of the chain $X^0$ on the same realization of
arrival times, service times and tie-breaking draws: $X$ started from
$(\gp,b)$ and $X'$ started from $(1,b)$. We claim that there is a stopping
time $\zeta$, equal to the first time $X$ is gap-free, such that for $t<\zeta$
\[
  X(t)=(\gp,\beta(t))\quad\text{and}\quad X'(t)=(1,\beta(t))
\]
for a common process $\beta$. Indeed, suppose the two chains are in
$(\gp,\beta)$ and $(1,\beta)$ and consider the next event.
\begin{itemize}
\item \emph{Service time of $B$.} Both move to the corresponding state with
  $(\beta-1)^+$ in place of $\beta$.
\item \emph{Service time of $A$.} Both move to $(0,\beta)$.
\item \emph{Arrival with $\beta\ge2$.} $X$ moves to $(2,\beta)$. In $X'$ the
  agent joins $A$ because $1<\beta$, so $X'$ moves to $(2,\beta)$.
\item \emph{Arrival with $\beta=1$.} $X$ moves to $(2,1)$. In $X'$ the agent is
  indifferent: if $u=A$ then $X'$ moves to $(2,1)$, and if $u=B$ then $X'$
  moves to $(1,2)$.
\item \emph{Arrival with $\beta=0$.} In $X$ the agent is indifferent, and in
  $X'$ she joins $B$. If $u=B$, the chains move to $(\gp,1)$ and $(1,1)$. If
  $u=A$, $X$ moves to $(2,0)$ and $X'$ moves to $(1,1)$.
\end{itemize}
This proves the claim, with $\zeta$ the first service time of $A$ or the first
arrival other than an arrival with $\beta=0$ and $u=B$. Moreover, at time
$\zeta$ either $X(\zeta)=X'(\zeta)$, or
\[
  E_1:\ \bigl(X(\zeta),X'(\zeta)\bigr)=\bigl((2,1),(1,2)\bigr),
  \qquad\text{or}\qquad
  E_0:\ \bigl(X(\zeta),X'(\zeta)\bigr)=\bigl((2,0),(1,1)\bigr).
\]

On $[0,\zeta)$ the two chains have the same numbers of agents in each queue, so
$r(X(s))=r(X'(s))$. By definition \eqref{eq:hgap},
$h(\gp,b)=\E[\int_0^\zeta r(X(s))ds+h(X(\zeta))]$. By
Lemma~\ref{lem:dynkin}, \eqref{eq:poisson} and optional stopping at
$\zeta\wedge t$,
$h(1,b)=\E[\int_0^{\zeta\wedge t}r(X'(s))ds+h(X'(\zeta\wedge t))]$; since
$|h(X'(\zeta\wedge t))|\le(1+b+N(\zeta))/\lambda$, which is integrable,
dominated convergence gives the same identity with $\zeta$ in place of
$\zeta\wedge t$. Subtracting and using $h(1,2)=-h(2,1)$ and $h(1,1)=0$,
\[
  \delta_b=\E\bigl[h(X(\zeta))-h(X'(\zeta))\bigr]
  =2h(2,1)\,\Pr_b(E_1)+h(2,0)\,\Pr_b(E_0).
\]
Both coefficients are positive by Lemma~\ref{lem:h}(a) and
$\Pr_b(E_1)+\Pr_b(E_0)\le1$, which gives the upper bound. For positivity: if
$b\ge1$, with positive probability the first $b-1$ events are service times of
$B$ and the next event is an arrival with $u=B$, which is $E_1$; if $b=0$, with
positive probability the first event is an arrival with $u=A$, which is $E_0$.

The final claim follows from Lemma~\ref{lem:h}(a) for gap-free states, and from
$|h(\gp,b)|\le|h(1,b)|+\bar\delta\le|1-b|/\lambda+\bar\delta$ for gap states.
\end{proof}

\begin{brem}
In both events $E_0$ and $E_1$ the gap has moved one agent from $B$ to $A$
relative to FCFS. The proof shows that this is the only effect a gap ever has,
and Lemma~\ref{lem:h}(a) shows that such a move is always valuable.
\end{brem}

\subsection{The advantage of the deviation}

For $b\ge0$ let
\[
  \mathrm{adv}(b)=\frac{1}{\lambda}\,(\cL_1h-\cL_0h)(0,b)
  =\E_u\,h(\mathrm{arr}_1((0,b),u))-\E_u\,h(\mathrm{arr}_0((0,b),u)).
\]

\begin{blem}\label{lem:adv}
$\mathrm{adv}(b)=\delta_b>0$ for $b\ge1$, and
$\mathrm{adv}(0)=-1/(1+\lambda)$.
\end{blem}

\begin{proof}
If $b\ge1$, then $\mathrm{arr}_1((0,b),u)=(\gp,b)$ and
$\mathrm{arr}_0((0,b),u)=(1,b)$, so $\mathrm{adv}(b)=h(\gp,b)-h(1,b)=\delta_b$.
If $b=0$, then $\mathrm{arr}_1((0,0),u)=(0,1)$, while
$\mathrm{arr}_0((0,0),u)$ is $(1,0)$ or $(0,1)$ with probability $1/2$ each, so
$\mathrm{adv}(0)=h(0,1)-\tfrac12(h(1,0)+h(0,1))=-h(1,0)$, and
$h(1,0)=1/(1+\lambda)$ by Lemma~\ref{lem:h}(b).
\end{proof}

\begin{bprop}[Performance difference]\label{prop:pd}
Under $\Reps$ with optimal routing, for every $T\ge0$,
\[
  \E\int_0^T r(X(s))\,ds
  =-\E\,h(X(T))+\lambda\,\E\int_0^T D(s)\,\mathrm{adv}(b(s))\,ds .
\]
\end{bprop}

\begin{proof}
By Lemma~\ref{lem:gap}, $h$ has at most linear growth, so
Lemma~\ref{lem:dynkin} applies to $f=h$. Since $D(s)=1$ only at states
$(0,b)$, we have $\cL_{D(s)}h(X(s))=\cL_0h(X(s))+\lambda D(s)\,
\mathrm{adv}(b(s))$, and $\cL_0h=-r$ by \eqref{eq:poisson}. Taking expectations
in Lemma~\ref{lem:dynkin} and using $h(X(0))=h(0,0)=0$ gives the result.
\end{proof}

Since $\mathrm{adv}(b)\ge0$ for $b\ge1$, Proposition~\ref{prop:pd} and
Lemma~\ref{lem:adv} give
\begin{equation}\label{eq:pdbound}
\begin{aligned}
  \E\int_0^T r(X(s))\,ds\ \ge\ {}&-\E|h(X(T))|
  +\lambda\,\delta_2\,\E\int_0^T D(s)\ind\{b(s)=2\}\,ds\\
  &-\frac{\lambda}{1+\lambda}\,\E\int_0^T D(s)\ind\{b(s)=0\}\,ds .
\end{aligned}
\end{equation}
It remains to show that $A$ is armed with $b=2$ much more often than it is
armed with $b=0$.

\subsection{Stability}

Let $V(x)=a(x)^2+b(x)^2$, and set $\theta=2-\lambda>0$ and $K=3\lambda+2$.

\begin{blem}\label{lem:drift}
For every $x\in\cX$ and every $d$ for which $\cL_d$ is defined at $x$,
\[
  \cL_dV(x)\le K-\theta\,n(x).
\]
Consequently, under $\Reps$ with optimal routing and for every $T>0$,
\[
  \E\,V(X(T))\le KT,\qquad \int_0^T\E\,n(X(s))\,ds\le\frac{KT}{\theta},
  \qquad \E\,n(X(T))\le\sqrt{2KT}.
\]
\end{blem}

\begin{proof}
A service time of $A$ changes $V$ by $-2a+1$ if $a>0$ and by $0$ otherwise, and
similarly for $B$; the two service terms sum to at most $-2n+2$. An arrival
that joins $A$ changes $V$ by $2a+1$, and one that joins $B$ by $2b+1$. We
check that the queue joined always holds at most $\min\{a,b\}+1$ agents. At
$(k,b)$ with $d=0$ the agent joins a queue of size $\min\{k,b\}$. At $(\gp,b)$
with $b\ge1$ she joins $A$, which has size $1=\min\{1,b\}$; at $(\gp,0)$ she
joins a queue of size at most $1=\min\{1,0\}+1$. At $(0,b)$ with $d=1$ she
joins $A$, of size $0$, if $b\ge1$, and $B$, of size $0$, if $b=0$. Hence the
arrival term is at most $\lambda(2\min\{a,b\}+3)\le\lambda(n+3)$, and
$\cL_dV\le\lambda(n+3)-2n+2=K-\theta n$.

Lemma~\ref{lem:dynkin} with $f=V$ and $X(0)=(0,0)$ gives
$0\le\E V(X(T))\le KT-\theta\int_0^T\E n(X(s))ds$, which yields the first two
bounds. The third follows from $n^2\le2V$ and Jensen's inequality.
\end{proof}

Let $\mathcal{R}$ denote the set of configurations in which both queues are
empty and $A$ is not armed.

\begin{blem}[Regeneration]\label{lem:regen}
Let $\varepsilon\in(0,1]$. Under $\Reps$ with optimal routing, the successive
entrance times $0=T_0<T_1<T_2<\cdots$ into $\mathcal{R}$ are almost surely
finite, the cycles of the process $(X,D,N_A,N_B)$ between them are i.i.d., and
$\E[T_1]<\infty$.
\end{blem}

\begin{proof}
\emph{Regeneration.} Suppose both queues are empty and $A$ is not armed at a
stopping time $T$. Until the next join to $A$, the quantities $j$ and $\ell$
do not change and $t-\ell(t)$ increases, so $A$ remains not armed; after the
next join, $D$ is determined by events after $T$. Hence the evolution of
$(X,D)$ after $T$ is the same measurable function of the arrival times, service
times and tie-breaking draws after $T$ as the evolution from time $0$ is of
those after $0$. By the strong Markov property of the Poisson processes, the
process after $T$ is independent of $\cF_T$ and has the same law as the
original process.

\emph{Finite mean.} Let $M=(K+1)/\theta$, $F=\{x\in\cX:n(x)<M\}$ and
$m_1=\lceil M\rceil$. By Lemma~\ref{lem:drift}, $\cL_dV\le-1$ outside $F$.
Starting from an arbitrary configuration with $X=x$ at a stopping time $\rho$,
let $\sigma=\inf\{t\ge\rho:X(t)\in F\}$. Lemma~\ref{lem:dynkin} applied to $V$
and optional stopping at $\sigma\wedge t$ give
$\E[\sigma\wedge t-\rho\mid\cF_\rho]\le V(x)$, hence
\begin{equation}\label{eq:hitF}
  \E[\sigma-\rho\mid\cF_\rho]\le V(X(\rho)).
\end{equation}
Next, consider the event $H$ that on $(\sigma,\sigma+1]$ there are no
arrivals, at least $m_1$ service times of $A$ and at least $m_1$ service times
of $B$. On $H$ both queues are empty at time $\sigma+1$, and since no agent has
joined $A$ during $(\sigma,\sigma+1]$ and $\varepsilon\le1$, $A$ is not armed
at $\sigma+1$. So the configuration at $\sigma+1$ lies in $\mathcal{R}$. Also
\begin{align*}
  \Pr(H\mid\cF_\sigma)&=p_0:=e^{-\lambda}\,\Pr(\Poi(1)\ge m_1)^2>0,\\
  \E\bigl[V(X(\sigma+1))\mid\cF_\sigma\bigr]&\le c_1:=\E\bigl[(m_1+\Poi(\lambda))^2\bigr],
\end{align*}
because $V\le n^2$ and $n(X(\sigma+1))\le n(X(\sigma))+N(\sigma,\sigma+1]$.

Now define rounds: $\rho_1=0$, and for $i\ge1$ let
$\sigma_i=\inf\{t\ge\rho_i:X(t)\in F\}$, let $H_i$ be the event $H$ for
$\sigma=\sigma_i$, and let $\rho_{i+1}=\sigma_i+1$. Let
$I=\min\{i:H_i\text{ occurs}\}$. The configuration lies in $\mathcal{R}$ at
time $\rho_{I+1}$. The event $\{I\ge i\}$ belongs to $\cF_{\rho_i}$, and
$\Pr(I\ge i)\le(1-p_0)^{i-1}$. By \eqref{eq:hitF},
\[
  \E[\rho_{I+1}]=\sum_{i\ge1}\E\bigl[(\rho_{i+1}-\rho_i)\ind\{I\ge i\}\bigr]
  \le\sum_{i\ge1}\E\bigl[(V(X(\rho_i))+1)\ind\{I\ge i\}\bigr].
\]
For $i\ge2$, since $\{I\ge i\}\subseteq\{I\ge i-1\}\in\cF_{\sigma_{i-1}}$,
\[
  \E\bigl[V(X(\rho_i))\ind\{I\ge i\}\bigr]
  \le\E\bigl[\ind\{I\ge i-1\}\,\E[V(X(\rho_i))\mid\cF_{\sigma_{i-1}}]\bigr]
  \le c_1(1-p_0)^{i-2}.
\]
Hence $\E[\rho_{I+1}]\le V(X(0))+(1+c_1)/p_0<\infty$ from any initial
configuration, and in particular from the configuration reached at the first
arrival after time $0$. Since the process leaves $\mathcal{R}$ at that arrival
and $T_1$ is at most the first subsequent time at which the configuration lies
in $\mathcal{R}$, we get $\E[T_1]\le1/\lambda+1+(1+c_1)/p_0<\infty$.
\end{proof}

\subsection{How often the rule is armed}

Call a join to $A$ a \emph{slot-$2$ join} if the agent is placed in slot $2$.
By Lemma~\ref{lem:states} this happens exactly when the state immediately
before is $(2,b)$, and by Lemma~\ref{lem:routing}(i) it requires $b\ge2$.

\begin{blem}[Armed against an empty $B$]\label{lem:bad}
For every $\varepsilon\in(0,1]$ and $T\ge0$,
\[
  \E\int_0^T D(s)\ind\{b(s)=0\}\,ds\le\frac{\lambda\,\varepsilon^{6}}{12}\,T .
\]
\end{blem}

\begin{proof}
Let $\vartheta_1<\vartheta_2<\cdots$ be the times of slot-$2$ joins. Suppose
$D(s)=1$ and $b(s)=0$, and let $\vartheta=\ell(s)$. Then $\vartheta$ is a
slot-$2$ join with $\vartheta<s<\vartheta+\varepsilon$. Immediately after
$\vartheta$, queue $A$ holds three agents and queue $B$ holds at least two.
Between $\vartheta$ and $s$ no agent joins $A$ and $A$ becomes empty, so there
are at least three service times of $A$ in $(\vartheta,s]$; and $b$ falls from
at least $2$ to $0$, so there are at least two service times of $B$ in
$(\vartheta,s]$. Therefore
\[
  \int_0^T D(s)\ind\{b(s)=0\}\,ds\le\varepsilon\sum_{i:\vartheta_i\le T}
  \ind\{\SvA(I_i)\ge3\}\,\ind\{\SvB(I_i)\ge2\},
  \qquad I_i=(\vartheta_i,\vartheta_i+\varepsilon].
\]
Each $\vartheta_i$ is a stopping time, so the two indicators are independent
of $\cF_{\vartheta_i}$ and of each other, with probabilities
$\Pr(\Poi(\varepsilon)\ge3)\le\varepsilon^3/6$ and
$\Pr(\Poi(\varepsilon)\ge2)\le\varepsilon^2/2$. The number of slot-$2$ joins up
to $T$ is at most $N(T)$, which has mean $\lambda T$. Taking expectations
gives the bound.
\end{proof}

\begin{blem}[Occupation of $(2,2)$]\label{lem:occupation}
Let $m_0=\lceil2K/\theta\rceil$ and
\[
  q=e^{-\lambda}\,\Pr(\Poi(1)\ge m_0)^2\cdot
  e^{-\lambda/2}\frac{(\lambda/2)^4}{4!}\cdot e^{-2}\cdot\frac14\cdot
  e^{-\lambda/2}>0 .
\]
For every $\varepsilon\in(0,1]$ and $T\ge2$, under $\Reps$ with optimal routing,
\[
  \int_0^T\Pr\bigl(X(s)=(2,2)\bigr)\,ds\ge\frac{q}{2}\,(T-2).
\]
\end{blem}

\begin{proof}
By Lemma~\ref{lem:drift} and Markov's inequality,
$\int_0^{T'}\Pr(n(X(s))\ge m_0)\,ds\le\frac{\theta}{2K}\cdot\frac{KT'}{\theta}
=T'/2$, so
\begin{equation}\label{eq:small}
  \int_0^{T'}\Pr\bigl(n(X(s))<m_0\bigr)\,ds\ge\frac{T'}{2}
  \qquad\text{for all }T'\ge0 .
\end{equation}
Fix $s\ge0$ and let $G_s$ be the event that all of the following hold:
\begin{enumerate}
\item on $(s,s+1]$ there are no arrivals, at least $m_0$ service times of $A$
  and at least $m_0$ service times of $B$;
\item on $(s+1,s+\tfrac32]$ there are exactly four arrivals, and the first and
  third of them have tie-breaking draw $A$;
\item on $(s+\tfrac32,s+2]$ there are no arrivals;
\item on $(s+1,s+2]$ there are no service times at either queue.
\end{enumerate}
Then $G_s$ is independent of $\cF_s$ and $\Pr(G_s)=q$. On
$G_s\cap\{n(X(s))<m_0\}$, both queues are empty at time $s+1$, and $A$ is not
armed at $s+1$ because no agent has joined $A$ for at least one unit of time
and $\varepsilon\le1$. The four arrivals then find, in order: state $(0,0)$,
where the agent is indifferent and joins $A$; state $(1,0)$, where the agent
joins $B$; state $(1,1)$, where the agent is indifferent and joins $A$; and
state $(2,1)$, where the agent joins $B$. Queue $A$ is not armed at any of
these arrivals, since it is either non-empty or its last join was not a
slot-$2$ join. Hence $X(u)=(2,2)$ for all $u\in[s+\tfrac32,s+2]$, and in
particular $\Pr(X(s+2)=(2,2))\ge q\,\Pr(n(X(s))<m_0)$. Integrating over
$s\in[0,T-2]$ and using \eqref{eq:small} gives the claim.
\end{proof}

\begin{blem}[Armed against $b=2$]\label{lem:good}
For every $\varepsilon\in(0,1]$ and $T\ge3$,
\[
  \E\int_0^T D(s)\ind\{b(s)=2\}\,ds\ge
  \frac{\lambda\,q\,\varepsilon^{4}}{384\,e^{4}}\,(T-3).
\]
\end{blem}

\begin{proof}
Let $\vartheta'_1<\vartheta'_2<\cdots$ be the times of slot-$2$ joins that
occur from state $(2,2)$; these are the arrivals that find state $(2,2)$ and
have tie-breaking draw $A$. For each $i$ let $G'_i$ be the event that
\[
  \SvA(\vartheta'_i,\vartheta'_i+\tfrac{\varepsilon}{2}]\ge3,\qquad
  N(\vartheta'_i,\vartheta'_i+\varepsilon]=0,\qquad
  \SvB(\vartheta'_i,\vartheta'_i+\varepsilon]=0 .
\]
Immediately after $\vartheta'_i$ the state is $(3,2)$. On $G'_i$, queue $A$ is
empty from some time before $\vartheta'_i+\varepsilon/2$ until
$\vartheta'_i+\varepsilon$, its last join is the slot-$2$ join at
$\vartheta'_i$, and $b=2$ throughout. Hence $D(s)=1$ and $b(s)=2$ for all
$s\in[\vartheta'_i+\varepsilon/2,\vartheta'_i+\varepsilon)$. On $G'_i$ there is
no arrival in $(\vartheta'_i,\vartheta'_i+\varepsilon]$, so these intervals are
disjoint for different $i$, and
\[
  \int_0^T D(s)\ind\{b(s)=2\}\,ds\ge\frac{\varepsilon}{2}
  \sum_{i:\vartheta'_i\le T-\varepsilon}\ind_{G'_i}.
\]
The event $G'_i$ is independent of $\cF_{\vartheta'_i}$, and since
$\varepsilon\le1$ and $\lambda<2$,
\[
  \Pr(G'_i)=\Pr\bigl(\Poi(\tfrac{\varepsilon}{2})\ge3\bigr)e^{-(\lambda+1)\varepsilon}
  \ge e^{-\varepsilon/2}\frac{(\varepsilon/2)^3}{6}\,e^{-(\lambda+1)\varepsilon}
  \ge\frac{\varepsilon^{3}}{48\,e^{4}} .
\]
By the compensation formula, the expected number of $\vartheta'_i$ in
$[0,T-\varepsilon]$ equals
$\frac{\lambda}{2}\int_0^{T-\varepsilon}\Pr(X(s)=(2,2))\,ds$, which by
Lemma~\ref{lem:occupation} is at least $\frac{\lambda q}{4}(T-3)$. Combining,
\[
  \E\int_0^T D(s)\ind\{b(s)=2\}\,ds\ge\frac{\varepsilon}{2}\cdot
  \frac{\varepsilon^{3}}{48\,e^{4}}\cdot\frac{\lambda q}{4}\,(T-3).\qedhere
\]
\end{proof}

\subsection{Proof of Theorem~\ref{thm:blind}}

Let $\varepsilon\in(0,1]$. Combining \eqref{eq:pdbound} with
Lemmas~\ref{lem:bad} and~\ref{lem:good}, for $T\ge3$,
\[
  \E\int_0^T r(X(s))\,ds\ \ge\ -\E|h(X(T))|
  +\lambda^{2}\varepsilon^{4}\Bigl(\frac{\delta_2\,q}{384\,e^{4}}\,(T-3)
  -\frac{\varepsilon^{2}}{12(1+\lambda)}\,T\Bigr).
\]
By Lemmas~\ref{lem:gap} and~\ref{lem:drift},
$\E|h(X(T))|\le\sqrt{2KT}/\lambda+\bar\delta$ and
$\E|a(X(T))-b(X(T))|\le\sqrt{2KT}$. Lemma~\ref{lem:busy} therefore gives
\begin{equation}\label{eq:liminf}
  \liminf_{T\to\infty}\frac{1}{T}\,\E\bigl[N_A(T)-N_B(T)\bigr]\ \ge\
  \lambda^{2}\varepsilon^{4}\Bigl(\kappa-\frac{\varepsilon^{2}}{12(1+\lambda)}\Bigr),
  \qquad\kappa=\frac{\delta_2\,q}{384\,e^{4}} .
\end{equation}

By Lemma~\ref{lem:regen}, $(N_A,N_B)$ is a cumulative process with respect to
the renewal sequence $(T_m)$, with $\E[T_1]<\infty$ and
$N_A(T_1)+N_B(T_1)=N(T_1)$, which has mean $\lambda\,\E[T_1]<\infty$. By the
renewal reward theorem (see, e.g., Asmussen, 2003, Chapter~VI),
\[
  \frac{N_q(t)}{t}\to\alpha_q:=\frac{\E[N_q(T_1)]}{\E[T_1]}
  \quad\text{almost surely, and}\quad
  \frac{\E[N_q(t)]}{t}\to\alpha_q,\qquad q\in\{A,B\}.
\]
Since $N_A(t)+N_B(t)=N(t)$, we have $\alpha_A+\alpha_B=\lambda$, and by
\eqref{eq:liminf},
\[
  \alpha_A-\alpha_B\ \ge\ \lambda^{2}\varepsilon^{4}
  \Bigl(\kappa-\frac{\varepsilon^{2}}{12(1+\lambda)}\Bigr).
\]
Hence
$\alpha_A\ge\frac{\lambda}{2}+\frac{\lambda^{2}\varepsilon^{4}}{2}
\bigl(\kappa-\frac{\varepsilon^{2}}{12(1+\lambda)}\bigr)$, which exceeds
$\lambda/2$ whenever
$\varepsilon<\bar\varepsilon=\min\{1,\sqrt{12(1+\lambda)\kappa}\}$. \qed

\subsection{Remarks}

\begin{brem}[Magnitude]
The bound in Theorem~\ref{thm:blind} is of order $\varepsilon^{4}$ and the
theorem makes no attempt to optimize it. Solving for the stationary
distribution numerically, in a variant of $\Reps$ in which the window of
length $\varepsilon$ is replaced by an independent exponential timer with mean
$\varepsilon$, gives a gain in $A$'s share of about $3\times10^{-7}$ at
$\lambda=1$ and $\varepsilon=1/5$, falling by a factor of about $16$ each time
$\varepsilon$ is halved, in line with the $\varepsilon^{4}$ rate. For
$\varepsilon\ge1$ the same variant loses share, so the requirement that
$\varepsilon$ be small is not an artifact of the proof.
\end{brem}

\begin{brem}[What the rule uses]
The rule $\Reps$ depends on the clock and is therefore not state-based in the
sense of Section~3, and it relies on preemption: the agent placed in slot $1$
of an empty queue is displaced by a later arrival placed in slot $0$. The
short window is what converts the first inference ($b\ge2$ at the slot-$2$
join) into near-certainty at the time of the deviation. The performance
difference identity in Proposition~\ref{prop:pd} holds for any blind rule that
deviates from FCFS placement only by offering slot $1$ when empty, with
$\E[D(s)\,\mathrm{adv}(b(s))]$ computed under that rule. A blind rule that
ignores the clock, or one with a non-negligible advantage, would deviate
whenever the posterior probability of $\{b=0\}$ given $A$'s own history is
small relative to the posterior mean of $\delta_b$; we do not pursue this here.
\end{brem}

\begin{brem}[Relation to the proof of Theorem~2]
The proof of Theorem~2 establishes pathwise dominance of the Ledger rule over
FCFS. That argument is not available here: when the deviation is made against
an empty queue $B$, the rule $\Reps$ loses an arrival that FCFS would have
attracted with probability $1/2$. The argument above instead compares the two
rules in expectation, through the value function $h$ of FCFS.
\end{brem}

\begin{thebibliography}{9}
\bibitem{asmussen} Asmussen, S. (2003). \emph{Applied Probability and Queues},
  2nd ed. Springer.
\bibitem{bremaud} Br\'emaud, P. (1981). \emph{Point Processes and Queues:
  Martingale Dynamics}. Springer.
\end{thebibliography}
% ============================= END APPENDIX BODY =============================

\end{document}
